 \documentclass[12pt]{article}
 % LaTeX file created by GUIDE version 46.2 on 04/30/26 at 10:10
 %%%% Usage: latex tree.tex, dvips tree.tex, ps2pdf tree.tex
 %%%% or: xelatex tree.tex
 \usepackage{pstricks}
 %%%% If using pdflatex, replace above line with \usepackage[pdf]{pstricks}
 %%%% and use: pdflatex -shell-escape tree.tex
 \usepackage{pstricks,pst-node,pst-tree}
 \usepackage{geometry}
 \usepackage{xcolor}
 \usepackage{lscape}
 \definecolor{wheat}{rgb}{0.96, 0.87, 0.7}
 \definecolor{lightgray}{rgb}{0.9, 0.9, 0.9}
 \definecolor{aqua}{rgb}{0.0, 1.0, 1.0}
 \definecolor{mauve}{rgb}{0.88, 0.69, 1.0}
 \definecolor{skyblue}{RGB}{86,180,233}
 \definecolor{vermillion}{RGB}{213,94,0}
 \definecolor{purple}{RGB}{204,121,167}
 \definecolor{bluishgreen}{RGB}{0,158,115}
 \pagestyle{empty}
 \begin{document}
 %\begin{landscape}
 \begin{center}
\psset{linecolor=black,tnsep=1pt,tndepth=0cm,tnheight=0cm,treesep=2.0cm,levelsep=50pt,radius=10pt,fillstyle=solid}
  \pstree[treemode=D]{\Tcircle{ 1 }~[tnpos=l]{\shortstack[r]{\texttt{\detokenize{X17p_status}}\\=\texttt{\detokenize{Loss}}}}
 }{
    \Tcircle[fillcolor=vermillion]{ 2 }~{\shortstack[c]{\emph{70}\\      3.24\\60+}} 
    \Tcircle[fillcolor=vermillion]{ 3 }~{\shortstack[c]{\emph{932}\\      0.88\\60+}} 
 }
 \end{center}
GUIDE v.46.2 0.250-SE
piecewise-constant proportional hazards regression tree
for \texttt{\detokenize{surv_5yrs}}.
 % Tree constructed from 1002 observations
 % (excluding observations with non-positive weight or missing values
 % in \texttt{D}, \texttt{T}, \texttt{R} or \texttt{Z} variables)
 % and pruned from a tree with 27 terminal nodes.
 % Maximum number of split levels is 15 and minimum node sample size is 2.
At each split, an observation goes to the left branch 
 if and only if the condition is satisfied.
Sample size (in \emph{italics}), relative hazard, and median survival time printed below nodes.
 Terminal nodes with median survival times above and below 60+ (median at root node) are painted \colorbox{yellow}{yellow} and \colorbox{vermillion}{vermillion} respectively.
 Second best split variable at root node is \texttt{\detokenize{X8p_status}}.
 %\end{landscape}
 \end{document}
